\section{From Mechanics to Fields}

In analytical mechanics, the motion of a system is described by generalized
coordinates \(q^i(t)\). A Lagrangian \(L(q,\dot{q},t)\) defines the action
\[
    S[q]=\int_{t_0}^{t_1} L(q,\dot{q},t)\,\dd t.
\]
The physical path is obtained by requiring the first variation of the action to
vanish.

Field theory uses the same principle, but the unknown is now a field
\(\phi(t,\vect{x})\). The role of the Lagrangian is played by a Lagrangian
density \(\mathcal{L}\), because the action must be integrated over space and
time.

\begin{remarkbox}
The word density matters. In mechanics the Lagrangian is integrated over time.
In field theory the Lagrangian density is integrated over spacetime volume.
\end{remarkbox}
